\section{Laurent cohomology of vector bundles}
\label{sec:laurent}

We retain the smooth threefold $V$, the cycle $D$, and the notation of
the preceding sections.  In particular,
\begin{equation}\label{eq:laurent-width-input}
 M_e=rp^e,\qquad \frac{d_e}{M_e}\longrightarrow
 a_*:=\frac{p}{p+1},\qquad c_2(V)\cdot D=0.
\end{equation}
The line bundle $\cO_V(M_eD)$ is trivial through width $d_e$.
Choose, once and for all, real numbers
\begin{equation}\label{eq:phase-constants}
 1-\frac2p<c_0<c<2a_*-1.
\end{equation}
Such choices exist since
\[
 2a_*-1-\left(1-\frac2p\right)=\frac{2}{p(p+1)}>0.
\]
As $p\ge5$, we also have $1/2<c_0<c<a_*$.
Put $u_e=\lfloor cM_e\rfloor$.
The lower bound on $c_0$ permits the use of only first and second
tensor powers in Lemma~\ref{lem:phase-bound}; the upper bound on $c$
leaves the precision required for matching jets in
Theorem~\ref{thm:laurent-vanishing}.

\subsection{Concentration of leading jets}

The first step restricts the positions at which a long jet of a
semistable bundle can have a nonzero leading term.

\begin{lemma}\label{lem:phase-bound}
Let $G$ be a vector bundle that is semistable of slope zero for the
degree $c_1(-)\cdot H\cdot D$.  There is a constant $A_G$ such that,
for every sufficiently large $e$,
\[
 g^0_{\lfloor c_0M_e\rfloor}(i;G)>0
 \quad\Longrightarrow\quad
 \dist(i,M_e\Z)\le A_G.
\]
\end{lemma}

\begin{proof}
Write $s=\lfloor c_0M_e\rfloor$.  For large $e$ we have $s\le d_e$.
A unit of $\cO(M_eD)|_{sD}$ translates the upper weight of a jet
by $M_e$.  We can therefore replace $i$ by a centered representative
$v$, with $|v|\le M_e/2$, without changing the existence of its
leading class.

Suppose that no uniform bound on $|v|$ exists.  Passing to a sequence,
we may assume that $v\to+\infty$ or $v\to-\infty$.
In the first case $s>v$ for large $e$.  The jet maps to a global
section of $G(vD)/G$ and retains its nonzero leading term.  Its pole
order tends to infinity, contrary to Proposition~\ref{prop:polar-bounds}.

In the second case choose a real number $\lambda$ such that
\begin{equation}\label{eq:phase-lambda}
 \frac{1-c_0}{2}<\lambda<\frac1p,
 \qquad p\lambda>\frac12.
\end{equation}
The interval is nonempty by \eqref{eq:phase-constants}.
Since the periods increase by the factor $p$, we can choose $j\le e$
with
\[
 \lambda\le\frac{|v|}{M_j}\le p\lambda.
\]
Here $j\to\infty$ as $|v|\to\infty$.  After a subsequence the
ratios tend to a number $y\in[\lambda,p\lambda]$.
There is an integer $m\in\{1,2\}$ for which
\begin{equation}\label{eq:phase-tensor-choice}
 1-c_0<my<1.
\end{equation}
Indeed, if $y>1-c_0$, take $m=1$; then $y\le p\lambda<1$.
Otherwise take $m=2$.  Its lower bound follows from
$2y\ge2\lambda>1-c_0$, and its upper bound follows from
$2y\le2(1-c_0)<4/p<1$.

Restrict the original jet to width $\lfloor c_0M_j\rfloor$, take
its $m$th tensor power, and multiply by a unit section of
$\cO(M_jD)$ to that width.  Tensoring keeps at least the same width:
in a difference of two tensor products, every summand contains a
factor from the lower lattice.  The resulting jet is a jet of
$G^{\otimes m}$ at upper weight
\[
 w=M_j+mv.
\]
By \eqref{eq:phase-tensor-choice}, $w\to+\infty$ and
$w<\lfloor c_0M_j\rfloor$ with a margin tending to infinity.
It therefore gives a global polar part of $G^{\otimes m}$.

This polar part has unbounded actual order.  At some component of
$D$ an original coefficient has pole ratio greater than $v-1$.
The corresponding pure tensor coefficient, after multiplication by
the unit, has pole ratio greater than $M_j+m(v-1)$.
This tends to infinity.  Proposition~\ref{prop:polar-bounds} applies
to both $G$ and $G\otimes G$, giving the required contradiction.
\end{proof}

\subsection{A vanishing criterion}

We use $\widehat E$ for the completion of a bundle $E$ along $D$,
and $E_{\cR}=\widehat E\otimes_{\widehat{\cO}}\cR$.
The finite-cover, bounded-pole interpretation of these sheaves is
the one used in the jet construction.

\begin{theorem}\label{thm:laurent-vanishing}
For the prepared data and periods fixed above, let $E$ be a vector
bundle on $V$ with the following properties:
\begin{enumerate}[label=\textup{(\roman*)}]
\item $\chi(P_i^1(E))=0$ for every integer $i$.
\item For $J=E,E^!$, the spaces $H^0(V,J(nD)/J)$ have dimensions
bounded independently of the positive integer $n$.
\item For $J=E,E^!$ there is a constant $A_J$ such that
\[
 g^0_{u_e}(i;J)>0\quad\Longrightarrow\quad
 \dist(i,M_e\Z)\le A_J
\]
for all sufficiently large $e$.
\end{enumerate}
Then global sections of $E_{\cR}$ form a finite-dimensional
$k$-vector space.  Every Laurent \v Cech $1$-cocycle on a finite
affine cover adapted to the bundles and the divisor has an actual
Laurent primitive.  The same assertion holds after refining such a
cover.
\end{theorem}

\begin{proof}
We separate the argument into the endpoint patterns, their drops,
the middle row, and convergence.

\medskip\noindent
\emph{Endpoint patterns.}
For $e$ large, $u_e\le d_e$, so the page at width $u_e$ is periodic
modulo $M_e$.  By (iii), the bottom row can be supported only at a
fixed bounded set of centered positions.  At each such position the
descending sequence of leading images stabilizes as the width tends
to infinity.  The stable-image statement of
Lemma~\ref{lem:page-calculus} realizes every vector in the stabilized
image by a global completed section at that cut.
Consequently, for all sufficiently large $e$ there is one fixed
bottom pattern, at positions $v$ with fixed multiplicities, repeated
by translations by $M_e$.  Duality gives a fixed top pattern at
positions $w$.  We enlarge a constant $A$ so that all these $v,w$
lie in $[-A,A]$.

This already proves finite-dimensionality of global Laurent
sections.  A nonzero section has a least upper cut $i$: its poles
are bounded, and the filtration is separated toward negative
infinity.  Its nonzero leading class belongs to $g^0_{u_e}(i;E)$
for every large $e$.  With $i$ fixed and $M_e\to\infty$, condition
(iii) forces $i$ to be an unshifted position of the fixed pattern.
There are only finitely many such cuts, and each leading quotient
is finite-dimensional.  These quotients give a finite filtration of
the space of global sections.  The same observation applies to $E^!$.

\medskip\noindent
\emph{Endpoint drops.}
By the middle-row recurrence, a bottom drop of length $s$ at
$v+jM_e$ removes middle-row rank at $v-s+jM_e$; a top drop at
$w+jM_e$ removes rank at $w+s+jM_e$.  To determine these losses
over the next period, we must control their lengths, with
multiplicities, independently of $j$ whenever $p\nmid j$.
Consider a bottom position $i=v+jM_e$.  If $p\mid j$, the next
period $M_{e+1}=pM_e$ retains exactly the same multiplicity there.
Monotonicity therefore permits no drop between pages $u_e$ and
$u_{e+1}$.  If $p\nmid j$, all of this multiplicity disappears
by page $u_{e+1}$.  We claim that its drops occur at lengths
\begin{equation}\label{eq:drop-range}
 d_e\le s\le M_e+C
\end{equation}
for a fixed $C$, and that their multiplicities at each length are
independent of $j$ among integers prime to $p$.

There is no drop before length $d_e$, since translation by a unit
of $\cO(M_eD)|_{d_eD}$ identifies the leading images with those
at the unshifted stable position $v$.
Choose local formal unit lifts $t_\lambda$ of this trivializing
section of $\cO(M_eD)|_{d_eD}$, and write on overlaps
\[
 \frac{t_\lambda}{t_\mu}=1+b_{\lambda\mu},
 \qquad b_{\lambda\mu}\in\cI_D^{d_e}\widehat{\cO}_V,
 \qquad \cI_D=\cO_V(-D).
\]
Fix a constant $C$ for the moment and a width
\[
 d_e<k\le M_e+C.
\]
For large $e$, $k<2d_e$ by \eqref{eq:laurent-width-input}.
Normalize local representatives $\sigma_\lambda$ of a width-$k$
jet at $v+jM_e$ by putting
$\alpha_\lambda=\sigma_\lambda/t_\lambda^j$.
These are local sections of $\widehat{E(vD)}$ and satisfy
\begin{equation}\label{eq:normalized-jet-overlap}
 \alpha_\mu=(1+jb_{\lambda\mu})\alpha_\lambda
 \pmod{\cI_D^k\widehat{E(vD)}}.
\end{equation}
Indeed $(1+b_{\lambda\mu})^j=1+jb_{\lambda\mu}$ to this
precision, for positive or negative $j$ alike, because $k<2d_e$.
In particular, the cochain $\alpha$ glues to width $d_e$.

We next match $\alpha$ to a single global completed section $\beta$
in that untwisted lattice through width $k-d_e$.
Every completed section and every formal lattice map restricts to the
actual algebraic finite jets of Equation~\eqref{eq:jet-quotient}:
on each affine intersection this is the identity
$\widehat A/(f^s)=A/(f^s)$.  These restrictions commute with the
\v Cech maps and with projection to leading layers.
Start with $\beta_{<0}=0$.  After matching the first $n$ depths,
the sum $\beta_{<n}$ of the completed global sections already chosen
satisfies
\[
 \alpha-\beta_{<n}\in C_{v-n}^0(E),\qquad
 \partial(\alpha-\beta_{<n})\in C_{v-d_e}^1(E).
\]
The second inclusion persists because each subtracted section is
global.  Thus, for $n<k-d_e$, the residual is a finite jet at upper
weight $v-n$ of width at least $d_e-n$.  The inequalities
\begin{equation}\label{eq:matching-margin}
 d_e-n\ge 2d_e-M_e-O(1)>u_e,
 \qquad k-d_e<M_e/2-O(1)
\end{equation}
hold with margins growing linearly in $M_e$, by $c<2a_*-1$ and
$a_*>1/2$.  Thus $v-n$ is centered, and condition (iii) says that
its leading image is either zero or one of the unshifted stabilized
images.  In the latter case lift it to a completed global section
and subtract.  A finite induction constructs the asserted $\beta$.

If $p\nmid j,j'$, put $q=j'/j$ in the prime field and set
\[
 \alpha'=q\alpha+(1-q)\beta.
\]
Since $\beta$ is global, Equation~\eqref{eq:normalized-jet-overlap}
gives
\[
 \alpha'_\mu-(1+j'b_{\lambda\mu})\alpha'_\lambda
 =j'(1-q)b_{\lambda\mu}(\alpha_\lambda-\beta)
 \pmod{\cI_D^k\widehat{E(vD)}}.
\]
The right side vanishes: its two factors have depths $d_e$ and
$k-d_e$.  Retwisting by $t_\lambda^{j'}$ therefore gives a
width-$k$ jet at $v+j'M_e$.  It has the same normalized leading
term, since $k-d_e\ge1$.  Also
$\alpha'-\beta=q(\alpha-\beta)$, so the same $\beta$ reverses
the construction with $q^{-1}$.  The normalized leading-image
subspaces are consequently equal at every width under consideration.
Their ranks, and hence their drops at each length, are independent of
$j$ prime to $p$.

For $j=1$, property (ii) kills the entire leading image by width
$M_e+C$ once $C>A$ is sufficiently large.  Indeed this width is
at least the upper weight $M_e+v$, so a surviving leading jet
would give a global polar part of order tending to infinity.
The increasing spaces in (ii) stabilize and forbid such polar
parts.  Since $u_{e+1}>M_e+C$ for large $e$, this proves
\eqref{eq:drop-range} and the complete drop rule.  The identical
argument for $E^!$, followed by the duality of
Lemma~\ref{lem:page-calculus}, gives the corresponding top-row rule.

\medskip\noindent
\emph{The middle row.}
We use the nonnegativity and Euler comparison from
Proposition~\ref{prop:scalar-widths}, now with a multiset of drop
positions.
Let $T_e$ be the multiset of the numbers $v-s$ for the bottom
drops and $w+s$ for the top drops.  Each is counted with its drop
multiplicity for one translating integer prime to $p$.
Its total size is the sum of the ranks of the two fixed endpoint
patterns.  Put $H_e(i)=g^1_{u_e}(i;E)$.  The middle-row recurrence
in Lemma~\ref{lem:page-calculus} gives
\begin{equation}\label{eq:middle-distribution}
\begin{split}
 H_{e+1}(i)=H_e(i)
 &-\#\{t\in T_e:t\equiv i\pmod{M_e}\}\\
 &+\#\{t\in T_e:t\equiv i\pmod{pM_e}\}.
\end{split}
\end{equation}
The last term restores the unshifted copies, which do not drop.
All members of $T_e$ lie in the two intervals
\begin{equation}\label{eq:drop-bands}
 [-M_e-O(1),-d_e+O(1)]
 \quad\hbox{and}\quad
 [d_e-O(1),M_e+O(1)].
\end{equation}
Each interval has length less than $M_e$ for large $e$.
For any fixed residue modulo $M_e$, these intervals can occupy
at most two of its lifts modulo $pM_e$.  Since $p\ge5$, some lift
is missed.  Nonnegativity of $H_{e+1}$ at that lift and periodicity
of $H_e$ show that
\[
 H_e(i)\ge\#\{t\in T_e:t\equiv i\pmod{M_e}\}.
\]
Summing over a period gives equality of the two sides, by the Euler
identity in Lemma~\ref{lem:page-calculus} and property (i).
The inequality is therefore equality at every residue.
Equation~\eqref{eq:middle-distribution} now says that $H_{e+1}$
is exactly the distribution of $T_e$ modulo $pM_e$.
For each fixed integer $i$, the bands \eqref{eq:drop-bands} and
$d_e\to\infty$ imply
\begin{equation}\label{eq:fixed-weight-vanishing}
 g^1_{u_e}(i;E)=0\qquad\text{for all sufficiently large }e.
\end{equation}

\medskip\noindent
\emph{Convergence of primitives.}
Fixed-weight page vanishing alone does not bound the poles of
successive primitives.  Fix an integer upper bound $l$ for a
Laurent cocycle.  We will bound the source $z=x+s$ of every nonzero
incoming differential of length $s$ to a middle-row weight $x\le l$
by a function $B_l(x)$ satisfying
\begin{equation}\label{eq:primitive-cut-bound}
 \sup_{x\le l}B_l(x)<\infty,
 \qquad B_l(x)\longrightarrow-\infty\quad(x\to-\infty).
\end{equation}

At a sufficiently late period, the drop rule gives
$z=v+jM_e$, $p\nmid j$, and $d_e\le s\le M_e+C$.
If $j\ge2$, then $x\ge M_e-A-C$, which exceeds $l$ for
all sufficiently large $e$.  If $j=1$, the target must lie in
the fixed finite set of integers in $[-A-C,l]$.
Choose one index $e_0$, depending on $l$, beyond the thresholds
for the endpoint and drop rules, so that $M_e-A-C>l$ for
$e\ge e_0$ and every middle rank in this finite set is zero at
page $u_{e_0}$.  Equation~\eqref{eq:fixed-weight-vanishing}
permits this simultaneous choice.  Middle ranks only decrease,
so no nonzero incoming image can have $j=1$ at any later page.

Put $S_l=u_{e_0}\ge1$.  Earlier lengths $s<S_l$ give
$z\le x+S_l$.  Every later length belongs to a period
$u_e\le s<u_{e+1}$ with $e\ge e_0$; the preceding exclusions
leave only $j\le-1$.  In that case $z\le-M_e+A$, and hence
$M_e\le-z+A$.  Thus $x=z-s\ge2z-A-C$, or
\[
 z\le\frac{x+A+C}{2}.
\]
The real-valued upper bound
\[
 B_l(x)=\max\left\{x+S_l,\frac{x+A+C}{2}\right\}
 \qquad(x\le l)
\]
therefore bounds every possible source and satisfies both assertions
of Equation~\eqref{eq:primitive-cut-bound}.

Let $\gamma$ be an actual Laurent $1$-cocycle of upper cut $x\le l$.
Its leading cochain belongs to every cycle space $Z_s^{x,1}$.
By \eqref{eq:fixed-weight-vanishing}, it eventually belongs to the
boundary space $U_s^{x,1}$.  More precisely, the page construction
identifies
\[
 U_{s+1}^{x,1}/U_s^{x,1}
 =\im\bigl(d_s:E_s^{x+s,0}\longrightarrow E_s^{x,1}\bigr),
\]
where $E_s^{i,a}=Z_s^{i,a}/U_s^{i,a}$.
Thus $U_s^{x,1}$ stops increasing once $x+s>B_l(x)$.
By its definition, each increase from a source cut $z=x+s$
has a representative primitive in $C_z^0(E)$; the initial space
$U_1^{x,1}$ uses $C_x^0(E)$.
Consequently the leading cochain of $\gamma$ has a primitive of
cut at most $\max\{x,B_l(x)\}$.
Subtract that primitive's coboundary.  The new cocycle has strictly
smaller upper cut.  Repeating at successive descending cuts gives
primitives whose upper cuts tend to negative infinity by
\eqref{eq:primitive-cut-bound}.  Their sum converges in a single
bounded-pole completed lattice.  The differential is continuous,
and the filtration is separated, so its coboundary is $\gamma$.
All arguments use finite covers by affine charts and their affine
intersections, and remain valid after refinement.  This proves the
theorem.
\end{proof}

\subsection{Cartier algebra on the Laurent charts}

We specify the module structures before applying the criterion to
exact forms.  Write $F$ for absolute Frobenius and
$\Omega=\Omega_V^1$.  Let $B^1,Z^1\subset F_*\Omega$ be the
sheaves of exact and closed algebraic one-forms, with the
Frobenius module structure.

\begin{lemma}\label{lem:laurent-cartier-algebra}
On a smooth affine chart with \'{e}tale coordinates $x_1,x_2,x_3$
and a local equation $f$ for $D$, put
$R=\widehat A[1/f]$.  The monomials
$x^I=x_1^{I_1}x_2^{I_2}x_3^{I_3}$, $0\le I_j<p$, form a basis
of $R$ over $R^p$.  Ordinary differentials and the Cartier
isomorphism extend to this ring.  The algebraic bundles $B^1,Z^1$
and their defining exact sequences extend to the corresponding
Laurent exact and closed forms, with their Frobenius module structure.
In particular,
\begin{equation}\label{eq:cartier-sequences}
 0\longrightarrow\cO_V\longrightarrow F_*\cO_V
   \longrightarrow B^1\longrightarrow0,
 \qquad
 0\longrightarrow B^1\longrightarrow Z^1
   \xrightarrow{C}\Omega\longrightarrow0
\end{equation}
are sequences of vector bundles and remain exact on Laurent charts.
The images and kernels in the Frobenius de Rham complex are locally
free, as are the relevant quotients.  Moreover,
\begin{equation}\label{eq:frobenius-duality}
 (F_*\cO_V)^\vee\otimes\omega_V\simeq F_*\omega_V.
\end{equation}
\end{lemma}

\begin{proof}
The algebraic calculation is the $p$-basis form of Cartier's theorem
\cite[Theorem~1]{cartier1957}; see also
\cite[Theorem~7.2]{katz1970}.  We first check that its finite basis
survives the completion and localization used here.
For an \'{e}tale algebra the relative Frobenius square is a base
change square.  One can check this because the relative Frobenius
is \'{e}tale and induces isomorphisms on geometric fibers.
The polynomial-coordinate monomial basis over $p$th powers
therefore gives the asserted basis for $A$ over $A^p$.
The field $k$ is perfect.

Complete the $p$-power subring along $f^p$ and use this finite basis.
The $f$-adic and $f^p$-adic topologies on $A$ coincide, and Frobenius
identifies the completion of $A^p$ with $(\widehat A)^p$.
The basis thus persists after completion and then after inverting
$f$; in the latter step, $1/f=f^{p-1}/f^p$.

Every $k$-derivation kills $p$th powers.  Since each element has a
unique finite expansion in the displayed basis with coefficients
in $R^p$, its differential is determined by ordinary differentiation
of those monomials.  Conversely the coordinate derivations are
defined by that rule.  Hence $\Omega^1_{R/k}$ is free on $dx_j$.
There is no additional continuity assumption on differentials.

In one variable, the de Rham complex decomposes into the monomials
of degrees less than $p$; its cohomology has representatives $1$
in degree zero and $x^{p-1}dx$ in degree one.  Tensoring the three
one-variable complexes proves local freeness of images, kernels,
and quotients, and proves the inverse Cartier rule
\[
 f\longmapsto f^p,\qquad
 df\longmapsto[f^{p-1}df].
\]
The rule is intrinsic.  Leibniz holds directly, while additivity
for differentials follows modulo exact forms by differentiating
the mixed terms of the integral polynomial
$((X+Y)^p-X^p-Y^p)/p$ and reducing modulo $p$.
Thus the local isomorphisms glue.

More explicitly, the scalar extension of the algebraic Frobenius
complex to a Laurent chart is the actual Laurent de Rham complex
with Frobenius action, through
\[
 R\otimes_A F_*\Omega_A^\bullet\longrightarrow
 F_*\Omega_R^\bullet,
 \qquad a\otimes\eta\longmapsto a^p\eta.
\]
The monomial and coordinate-wedge bases make this an isomorphism.
This proves the assertions about \eqref{eq:cartier-sequences}.
Finally, the pairing obtained by Cartier on top forms is perfect:
in coordinates it extracts the coefficient of
$x_1^{p-1}x_2^{p-1}x_3^{p-1}dx_1\wedge dx_2\wedge dx_3$.
It gives \eqref{eq:frobenius-duality} and is intrinsic by the same
Cartier rule.
\end{proof}

\subsection{Application to one-forms and exact forms}

\begin{corollary}\label{cor:form-vanishing}
Every Laurent \v Cech $1$-cocycle with values in $\Omega$, $B^1$,
or $Z^1$ has a Laurent primitive.  Global Laurent one-forms form
a finite-dimensional $k$-vector space.
\end{corollary}

\begin{proof}
For $E=\Omega$, both $E$ and $E^!=T_V\otimes\omega_V$ are
semistable of slope zero by
Proposition~\ref{prop:strong-semistability} and the canonical
intersection identities.  Proposition~\ref{prop:polar-bounds} and
Lemma~\ref{lem:phase-bound} give conditions (ii) and (iii) of
Theorem~\ref{thm:laurent-vanishing}; passing from width
$\lfloor c_0M_e\rfloor$ to $u_e$ can only decrease the leading
images.  Hirzebruch--Riemann--Roch gives condition (i).
Indeed the degree-two terms in $\ch(\Omega)\td(T_V)$ are linear
combinations of $K_V^2$ and $c_2(V)$, both killed by intersection
with $D$, and every term involving $D^2$ is also zero.

For $E=B^1$, the first sequence in
\eqref{eq:cartier-sequences} gives condition (i).  By projection
formula, a single cell for $F_*\cO_V$ has the Euler characteristic
of $p$ successive scalar cells; each is zero by
Proposition~\ref{prop:scalar-widths}.  Frobenius changes the
$k$-module action but not these finite dimensions.

There are subbundle inclusions with locally free quotients
\begin{equation}\label{eq:exact-form-inclusions}
 B^1\longhookrightarrowactual F_*\Omega,
 \qquad
 (B^1)^!\longhookrightarrowactual F_*\omega_V.
\end{equation}
The first follows from Lemma~\ref{lem:laurent-cartier-algebra};
the second follows by dualizing the first sequence in
\eqref{eq:cartier-sequences} and using
\eqref{eq:frobenius-duality}.
For a subbundle $E\subset F_*G$, the injection remains injective
on every quotient $E(nD)/E$.  Projection formula embeds that
quotient in $F_*(G(pnD)/G)$.  The polar bounds for $G=\Omega$
and $G=\omega_V$ therefore give condition (ii) for both bundles
in \eqref{eq:exact-form-inclusions}.
All these bounds concern fixed bundles and their first Frobenius
pushforwards.  The growing index $e$ below changes the jet period, not
the bundle in the polar-part estimate.

The same subbundle property gives condition (iii), with a small
precision margin that is useful to make explicit.  A width-$u_e$
jet of $E\subset F_*G$ with a nonzero leading term maps to a
width-$pu_e$ jet of $G$ at upper weight $pi$.  Its leading term
is nonzero in the top $p$ ordinary layers.  Choose its first
nonzero ordinary layer, at a weight $a$ with $|a-pi|<p$.
It has width at least $pu_e-p$.  Since $c>c_0$ and
$M_{e+1}=pM_e$, for all sufficiently large $e$ we have
\[
 pu_e-p\ge\lfloor c_0M_{e+1}\rfloor.
\]
Lemma~\ref{lem:phase-bound} at period $M_{e+1}$ implies that $a$,
and hence $pi$, is at bounded distance from $M_{e+1}\Z$.
Dividing by $p$ gives the required bounded distance from $M_e\Z$.

Theorem~\ref{thm:laurent-vanishing} now applies to $\Omega$ and
$B^1$, including finite-dimensionality for the former.
For a $Z^1$ cocycle, first solve its image in $\Omega$ and lift
the resulting local primitives through Cartier.  After subtraction
the remaining cocycle lies in $B^1$ and can also be solved.
These lifts exist on affine Laurent charts by
Lemma~\ref{lem:laurent-cartier-algebra}; refining the finite cover
if necessary preserves every preceding argument.  This proves
the assertion for $Z^1$.
\end{proof}

The next section uses both conclusions: cocycle vanishing supplies
compatible closed connection forms, and finite-dimensionality permits
a global adjustment that makes those forms Cartier-fixed.
